\begin{afterword}
\summaryhead

The post warns against drawing from the anti-zombie argument the rule that whatever cannot be seen does not exist. A photon sent into intergalactic space will, because the expansion of the universe is accelerating, pass beyond any possible interaction with us; yet it would be silly to think that it then vanishes. Why is this not a licence to believe in any invisible thing?

The first answer is that Occam's razor, in its Bayesian forms, Solomonoff induction and Minimum Message Length, charges for the length of a model's laws, not for the number of things the laws govern. Extra galaxies cost nothing; a law that makes the photon vanish is costly. The post then refines this into a distinction. We assign probability to simple, testable laws and believe what they logically imply: the ``implied invisible.'' A specific unseen event or law with no evidence of its own, the ``additional invisible,'' keeps only its prior probability. An added paragraph asks whether it is worth sending colonists to a place from which no message could ever return.

\responsehead

The physics is right. In the standard cosmological model the universe has an event horizon, and a photon sent out now passes beyond it after a finite time; after that, nothing from the photon can reach us. The added case of the colony ship holds for the same reason, since the horizon shrinks in comoving terms while the ship travels.

The formal argument is correct for the formalizations it names, and it states its own limit: extra quarks cost bits if the model's predictions depend on their exact positions. One sentence goes further than the argument. ``Occam's Razor doesn't care about RAM'' is true of Solomonoff induction and Minimum Message Length, but not of every formalization proposed. Schmidhuber's speed prior, published in 2002, counts running time, because Solomonoff's measure gives high probability to data that are extremely hard to compute.

The distinction between the implied and the additional invisible refines the simplicity answer rather than replacing it. The equations with a vanishing clause fit all the evidence as well as the equations without it; they are rejected because the extra law costs prior probability, as the post itself says. What the distinction adds is a rule about where probability attaches: to laws, and to unseen things only through what the laws imply. The rule is offered with hedges (``does seem,'' ``I can't think of a case''). A similar point had been made before; Tegmark, whose essay on parallel universes the author has cited, wrote in 2003 that ``Containing unobservable entities does clearly not per se make a theory non-testable.''

Two historical remarks are loose, and neither affects the argument. The Latin line is not Ockham's; the phrasing comes from John Punch in 1639. The story that Occam's razor was invoked against other galaxies is offered from memory, with the post's own admission that it could not find a source, and I could not find one either.

The principle carries weight later in the book: ``Decoherence is Simple'' calls itself ``a restatement of the points in Belief in the Implied Invisible, as they apply to quantum physics.''

In short: a sound argument, with correct physics, that belief in tested laws carries belief in what they imply beyond sight; its distinction refines the simplicity answer rather than replacing it.
\end{afterword}
